\documentclass[11pt]{article}
\usepackage[margin=1in]{geometry}
\usepackage{amsmath, amssymb, amsthm}
\usepackage{enumitem}
\usepackage{xcolor}
\usepackage{tcolorbox}
\tcbuselibrary{skins, breakable}

% Define color boxes for different framework sections
\newtcolorbox{problemconfigbox}{
    colback=blue!5, colframe=blue!40!black, title=PROBLEM CONFIGURATION ANALYSIS,
    fonttitle=\bfseries, breakable
}

\newtcolorbox{strategicbox}{
    colback=pink!5, colframe=pink!40!black, title=STRATEGIC FRAMEWORK APPLICATION,
    fonttitle=\bfseries, breakable
}

\newtcolorbox{tacticalbox}{
    colback=green!5, colframe=green!40!black, title=TACTICAL METHODOLOGY SELECTION,
    fonttitle=\bfseries, breakable
}

\newtcolorbox{conversionbox}{
    colback=yellow!5, colframe=yellow!40!black, title=CONVERSION TECHNIQUES APPLICATION,
    fonttitle=\bfseries, breakable
}

\newtcolorbox{verificationbox}{
    colback=red!5, colframe=red!40!black, title=VERIFICATION STRATEGY DESIGN,
    fonttitle=\bfseries, breakable
}

\newtcolorbox{roadmapbox}{
    colback=cyan!5, colframe=cyan!40!black, title=SOLUTION ROADMAP,
    fonttitle=\bfseries, breakable
}

\newtcolorbox{competitionbox}{
    colback=purple!5, colframe=purple!40!black, title=COMPETITION STRATEGY INTEGRATION,
    fonttitle=\bfseries, breakable
}

\newtcolorbox{keyinsightbox}{
    colback=orange!5, colframe=orange!40!black, title=KEY INSIGHT,
    fonttitle=\bfseries, breakable
}

\newtcolorbox{warningbox}{
    colback=red!10, colframe=red!60!black, title=WARNING,
    fonttitle=\bfseries, breakable
}

\newtcolorbox{tipbox}{
    colback=green!10, colframe=green!60!black, title=PRO TIP,
    fonttitle=\bfseries, breakable
}

\title{\textbf{2016 AMC 12B Problem 24: Strategic Analysis}}
\author{Using AMC12/COMC Problem-Solving Framework}
\date{}

\begin{document}
\maketitle

\section*{Problem Statement}
There are exactly $77,000$ ordered quadruplets $(a, b, c, d)$ such that $\gcd(a, b, c, d) = 77$ and $\operatorname{lcm}(a, b, c, d) = n$. What is the smallest possible value for $n$?

\textbf{Answer Choices:} 
$\textbf{(A)}\ 13,860\qquad\textbf{(B)}\ 20,790\qquad\textbf{(C)}\ 21,560 \qquad\textbf{(D)}\ 27,720 \qquad\textbf{(E)}\ 41,580$

\begin{problemconfigbox}
\subsection*{A. Problem Classification}
\begin{itemize}[leftmargin=*]
    \item \textbf{Problem Type}: To Find - determine the minimum value of $n$ given counting constraints
    \item \textbf{Difficulty Band}: Late-band (Problem 24 of 25) - requires advanced number theory and counting
    \item \textbf{Competition Context}: AMC12 (75 min, 25 problems) - allocated time $\sim$5-8 minutes
    \item \textbf{Mathematical Domain}: Number Theory (GCD/LCM properties) + Combinatorics (counting with constraints)
\end{itemize}

\subsection*{B. Problem Structure Identification}
\begin{itemize}[leftmargin=*]
    \item \textbf{Given Information}: 
    \begin{itemize}
        \item Exactly $77,000$ ordered quadruplets $(a,b,c,d)$
        \item $\gcd(a,b,c,d) = 77 = 7 \times 11$
        \item $\operatorname{lcm}(a,b,c,d) = n$ (unknown)
    \end{itemize}
    \item \textbf{Constraints}: 
    \begin{itemize}
        \item Must minimize $n$
        \item Count must be exactly $77,000$
        \item All numbers are positive integers
    \end{itemize}
    \item \textbf{Objective}: Find minimum value of $n$ satisfying the counting condition
    \item \textbf{Hidden Assumptions}: 
    \begin{itemize}
        \item The relationship between prime factorizations of GCD and LCM
        \item The counting must factor nicely to $77,000 = 77 \times 1000 = 7 \times 11 \times 8 \times 125$
    \end{itemize}
    \item \textbf{Answer Choice Leverage}: All answers are multiples of $77$; we need to find which prime factorization works
\end{itemize}

\subsection*{C. Strategic Orientation}
\begin{itemize}[leftmargin=*]
    \item \textbf{Hypothesis}: After dividing by $\gcd = 77$, we get relatively prime numbers whose LCM structure determines the count
    \item \textbf{Conclusion}: Minimize $n$ by assigning larger exponents to smaller primes
    \item \textbf{Notation}: Let $A = a/77, B = b/77, C = c/77, D = d/77$ with $\gcd(A,B,C,D) = 1$ and $N = n/77 = \operatorname{lcm}(A,B,C,D)$
    \item \textbf{Intuitive Approach}: Count quadruplets based on exponents in prime factorization
    \item \textbf{Cross-Domain Potential}: This is fundamentally a combinatorics problem disguised as number theory
\end{itemize}
\end{problemconfigbox}

\begin{strategicbox}
\subsection*{A. Psychological Strategies}
\begin{itemize}[leftmargin=*]
    \item \textbf{Mental Toughness}: This is P24 - expect complexity. Don't panic if first approach doesn't immediately work. The key is systematic analysis of the counting function.
    \item \textbf{Creativity Cultivation}: The key insight is recognizing that we need to factor $77,000$ as a product of counting functions $f(k)$
    \item \textbf{Iterative Refinement}: Start with small cases to understand the counting pattern, then generalize
\end{itemize}

\subsection*{B. Investigation Strategies}
\begin{itemize}[leftmargin=*]
    \item \textbf{Startup Orientation}: Recognize this as a GCD/LCM problem requiring normalization (divide by GCD)
    \item \textbf{Get Your Hands Dirty}: Compute $f(1), f(2), f(3)$ to find the pattern: $f(k) = 12k^2 + 2$
    \item \textbf{Penultimate Step}: Once we have $77000 = f(k_1) \times f(k_2) \times \cdots$, assign exponents optimally
    \item \textbf{Wishful Thinking}: "I wish the counting function had a simple closed form" - this guides us to derive $f(k)$
    \item \textbf{Make It Easier}: Consider the case of pairs $(a,b)$ first to understand the counting mechanism
    \item \textbf{Conjecture via Small Cases}: 
    \begin{itemize}
        \item $f(1) = 14$ (quadruplets with max exponent 1, min exponent 0)
        \item $f(2) = 50$
        \item Pattern: $f(k) = 14 + 12(k^2-1) = 12k^2 + 2$
    \end{itemize}
    \item \textbf{Visualization}: Think of exponents in a grid where each prime power contributes independently
\end{itemize}

\subsection*{C. Late-Band Strategic Playbook (P17-25)}
\begin{itemize}[leftmargin=*]
    \item \textbf{Answer-Choice Leverage}: Factor each answer choice: $27720 = 77 \times 360 = 77 \times 2^3 \times 3^2 \times 5$
    \item \textbf{Make It Easier}: Replace quadruplets with pairs to understand basic counting, then extend
    \item \textbf{Penultimate Step}: The final step is computing $N \times 77$ after finding optimal exponent assignment
    \item \textbf{Wishful Thinking}: Assume $77000$ factors nicely into products of $f(k)$ values
    \item \textbf{Conjecture via Small Cases}: $14 \times 50 \times 110 = 77000$ suggests exponents $(1,2,3)$
\end{itemize}

\subsection*{D. Argument Construction Strategy}
\begin{itemize}[leftmargin=*]
    \item \textbf{Direct Proof}: Derive the counting function $f(k)$ from first principles, then factor $77000$
    \item \textbf{Proof by Cases}: Handle each prime power independently in the factorization
    \item \textbf{Contradiction Path}: Not applicable here
    \item \textbf{Constructive Proof}: Show that $N = 2^3 \times 3^2 \times 5^1 = 360$ achieves the count
\end{itemize}
\end{strategicbox}

\begin{tacticalbox}
\subsection*{A. Core Mathematical Tactics}
\begin{itemize}[leftmargin=*]
    \item \textbf{GCD/LCM Normalization}: Divide all numbers by $\gcd(a,b,c,d) = 77$ to get $\gcd(A,B,C,D) = 1$
    \item \textbf{Prime Factorization}: Express $N$ and all variables in terms of prime powers
    \item \textbf{Exponent Analysis}: For each prime $p$, the exponents must include both maximum and minimum (0 and $k$)
    \item \textbf{Counting Principle}: Count ordered quadruplets by considering exponent distributions
    \item \textbf{Optimization Strategy}: Assign larger exponents to smaller primes to minimize $N$
\end{itemize}

\subsection*{B. Domain-Specific Approaches}
\begin{itemize}[leftmargin=*]
    \item \textbf{Number Theory}: Use fundamental theorem of arithmetic and GCD/LCM properties
    \item \textbf{Combinatorics}: Develop counting function for exponent distributions
    \item \textbf{Algebraic Manipulation}: Derive closed form $f(k) = 12k^2 + 2$ using complementary counting
\end{itemize}

\subsection*{C. Crossover Tactics}
\begin{itemize}[leftmargin=*]
    \item \textbf{Generating Functions}: Not directly applicable but thinking about products of terms helps
    \item \textbf{Inclusion-Exclusion}: Used in deriving $f(k) = (k+1)^4 - k^4 - k^4 + (k-1)^4$
\end{itemize}
\end{tacticalbox}

\begin{conversionbox}
\subsection*{A. Recognition Index - High Probability Techniques}
\begin{itemize}[leftmargin=*]
    \item \textbf{Technique \#1: GCD-LCM Normalization}
    \begin{itemize}
        \item \textit{Why it applies}: Problem explicitly involves both GCD and LCM
        \item \textit{Trigger}: When you see $\gcd(a,b,c,d) = g$ and $\operatorname{lcm}(a,b,c,d) = n$
        \item \textit{Application}: Set $A = a/77$, etc., to get $\gcd(A,B,C,D) = 1$
    \end{itemize}
    
    \item \textbf{Technique \#4: Prime Factorization Decomposition}
    \begin{itemize}
        \item \textit{Why it applies}: LCM/GCD inherently depend on prime factorizations
        \item \textit{Trigger}: Number theory problem with divisibility
        \item \textit{Application}: Write $N = 2^{k_2} \times 3^{k_3} \times 5^{k_5} \times \cdots$
    \end{itemize}
    
    \item \textbf{Technique \#8: Strategic Counting with Constraints}
    \begin{itemize}
        \item \textit{Why it applies}: Must count exactly $77,000$ quadruplets
        \item \textit{Trigger}: "How many" with specific constraint value
        \item \textit{Application}: Develop function $f(k)$ for each prime power
    \end{itemize}
    
    \item \textbf{Technique \#11: Pattern Recognition from Small Cases}
    \begin{itemize}
        \item \textit{Why it applies}: Complex counting benefits from concrete examples
        \item \textit{Trigger}: When formula isn't immediately obvious
        \item \textit{Application}: Compute $f(1) = 14, f(2) = 50, f(3) = 110$ to find pattern
    \end{itemize}
\end{itemize}

\subsection*{B. Technique Integration Strategy}
\begin{itemize}[leftmargin=*]
    \item \textbf{Primary Technique}: GCD-LCM Normalization + Prime Factorization (Techniques \#1, \#4)
    \item \textbf{Secondary Support}: Pattern recognition (Technique \#11) to derive $f(k) = 12k^2 + 2$
    \item \textbf{Final Step}: Optimization (assign larger exponents to smaller primes)
    \item \textbf{Synergy}: The techniques work sequentially - normalize first, then factor, then count, then optimize
\end{itemize}

\subsection*{C. Advanced Conversion Analysis}
\begin{itemize}[leftmargin=*]
    \item \textbf{Complementary Counting Formula}:
    \[f(k) = (k+1)^4 - k^4 - k^4 + (k-1)^4 = 12k^2 + 2\]
    \item \textbf{Factorization}: $77000 = 8 \times 125 \times 77 = 2^3 \times 5^3 \times 7 \times 11$
    \item \textbf{Key Realization}: $14 \times 50 \times 110 = 77000$, where:
    \begin{itemize}
        \item $f(1) = 14 = 12(1)^2 + 2$
        \item $f(2) = 50 = 12(4) + 2$
        \item $f(3) = 110 = 12(9) + 2$
    \end{itemize}
\end{itemize}
\end{conversionbox}

\begin{verificationbox}
\subsection*{A. Verification Level Selection}
Using the decision tree for a late-band problem with high complexity:
\begin{itemize}[leftmargin=*]
    \item \textbf{Problem Band}: Late (P24) → Minimum Level 4
    \item \textbf{Confidence}: Medium (70-80\%) → Upgrade to Level 5
    \item \textbf{Selected Level}: Level 5 - Thorough Verification
\end{itemize}

\subsection*{B. Primary Verification Methods}
\begin{itemize}[leftmargin=*]
    \item \textbf{Forward Check}: Verify $f(1) \times f(2) \times f(3) = 14 \times 50 \times 110$:
    \begin{align*}
        14 \times 50 &= 700 \\
        700 \times 110 &= 77,000 \quad \checkmark
    \end{align*}
    
    \item \textbf{Formula Verification}: Check $f(k) = 12k^2 + 2$ for $k=1,2,3$:
    \begin{align*}
        f(1) &= 12(1) + 2 = 14 \quad \checkmark \\
        f(2) &= 12(4) + 2 = 50 \quad \checkmark \\
        f(3) &= 12(9) + 2 = 110 \quad \checkmark
    \end{align*}
    
    \item \textbf{Answer Computation}: $N = 2^3 \times 3^2 \times 5^1 = 8 \times 9 \times 5 = 360$
    \[n = 77 \times 360 = 27,720 \quad \text{(Answer D)}\]
\end{itemize}

\subsection*{C. Specialized Verification}
\begin{itemize}[leftmargin=*]
    \item \textbf{Optimality Check}: Verify we assigned exponents correctly:
    \begin{itemize}
        \item Largest exponent (3) to smallest prime (2): $2^3$
        \item Middle exponent (2) to middle prime (3): $3^2$
        \item Smallest exponent (1) to largest prime (5): $5^1$
        \item This minimizes $N = 360$
    \end{itemize}
    
    \item \textbf{Alternative Assignment Check}:
    \begin{itemize}
        \item If we used $2^1 \times 3^2 \times 5^3 = 2 \times 9 \times 125 = 2250$ (larger)
        \item If we used $2^2 \times 3^3 \times 5^1 = 4 \times 27 \times 5 = 540$ (larger)
        \item Our assignment gives minimum: $360 < 540 < 2250$ \quad \checkmark
    \end{itemize}
\end{itemize}

\subsection*{D. Confidence Calibration}
\begin{itemize}[leftmargin=*]
    \item \textbf{Initial Confidence}: 75\% (derived formula and found factorization)
    \item \textbf{After Verification}: 90\% (confirmed counting and optimality)
    \item \textbf{Remaining Uncertainty}: 10\% (potential arithmetic errors in $f(k)$ derivation)
\end{itemize}
\end{verificationbox}

\begin{roadmapbox}
\subsection*{A. Step-by-Step Solution Plan}

\textbf{Step 1: Normalize by GCD (2 minutes)}
\begin{itemize}[leftmargin=*]
    \item Let $A = a/77, B = b/77, C = c/77, D = d/77$
    \item Then $\gcd(A,B,C,D) = 1$ and $N = n/77 = \operatorname{lcm}(A,B,C,D)$
    \item The count of $(a,b,c,d)$ equals the count of $(A,B,C,D)$
    \item \textit{Checkpoint}: Do we understand why normalization preserves count? $\checkmark$
\end{itemize}

\textbf{Step 2: Analyze Exponent Structure (2 minutes)}
\begin{itemize}[leftmargin=*]
    \item Write $N = 2^{k_2} \times 3^{k_3} \times 5^{k_5} \times \cdots$
    \item For each prime $p^k$ in factorization of $N$:
    \begin{itemize}
        \item At least one of $(A,B,C,D)$ has exponent $k$ (for LCM)
        \item At least one has exponent $0$ (for GCD = 1)
        \item All have exponents between $0$ and $k$
    \end{itemize}
    \item \textit{Checkpoint}: Clear on constraint structure? $\checkmark$
\end{itemize}

\textbf{Step 3: Derive Counting Function (3 minutes)}
\begin{itemize}[leftmargin=*]
    \item Define $f(k)$ = number of ordered quadruplets $(m_1,m_2,m_3,m_4)$ where:
    \begin{itemize}
        \item $0 \le m_i \le k$ for all $i$
        \item $\max(m_i) = k$ and $\min(m_i) = 0$
    \end{itemize}
    \item Use complementary counting:
    \[f(k) = (k+1)^4 - k^4 - k^4 + (k-1)^4\]
    \item Simplify to get $f(k) = 12k^2 + 2$
    \item \textit{Checkpoint}: Verified formula with small cases? $\checkmark$
\end{itemize}

\textbf{Step 4: Factor the Count (2 minutes)}
\begin{itemize}[leftmargin=*]
    \item We need $f(k_2) \times f(k_3) \times f(k_5) \times \cdots = 77,000$
    \item Compute: $f(1) = 14, f(2) = 50, f(3) = 110, f(4) = 194, \ldots$
    \item Find: $14 \times 50 \times 110 = 77,000$ \quad $\checkmark$
    \item Therefore exponents are $(1, 2, 3)$
    \item \textit{Checkpoint}: Does factorization work? $\checkmark$
\end{itemize}

\textbf{Step 5: Optimize Prime Assignment (1 minute)}
\begin{itemize}[leftmargin=*]
    \item To minimize $N$, assign larger exponents to smaller primes:
    \begin{itemize}
        \item Exponent 3 → prime 2: $2^3 = 8$
        \item Exponent 2 → prime 3: $3^2 = 9$
        \item Exponent 1 → prime 5: $5^1 = 5$
    \end{itemize}
    \item Calculate: $N = 8 \times 9 \times 5 = 360$
    \item Final answer: $n = 77 \times 360 = 27,720$
    \item \textit{Checkpoint}: Is this the minimum? $\checkmark$
\end{itemize}

\subsection*{B. Alternative Approaches}
\begin{itemize}[leftmargin=*]
    \item \textbf{Backup Method 1}: If pattern not clear, compute more $f(k)$ values systematically
    \item \textbf{Backup Method 2}: Work backwards from answer choices to verify which gives count $77,000$
    \item \textbf{Pivot Indicator}: If after 3 minutes no pattern emerges, switch to answer verification
\end{itemize}

\subsection*{C. Critical Pitfalls Guide}
\begin{itemize}[leftmargin=*]
    \item \textbf{Pitfall 1}: Forgetting to multiply back by 77 at the end
    \begin{itemize}
        \item \textit{Prevention}: Write $n = 77 \times N$ prominently
        \item \textit{Detection}: Answer should be among choices
    \end{itemize}
    
    \item \textbf{Pitfall 2}: Counting function off by constant
    \begin{itemize}
        \item \textit{Prevention}: Verify $f(1)$ by hand-counting
        \item \textit{Detection}: $14 \times 50 \times 110$ should equal $77,000$
    \end{itemize}
    
    \item \textbf{Pitfall 3}: Wrong exponent assignment (not minimizing)
    \begin{itemize}
        \item \textit{Prevention}: Remember larger exponents on smaller primes
        \item \textit{Detection}: Check all permutations give larger values
    \end{itemize}
    
    \item \textbf{Pitfall 4}: Arithmetic errors in computing $f(k) = 12k^2 + 2$
    \begin{itemize}
        \item \textit{Prevention}: Double-check with formula and verify with small cases
        \item \textit{Detection}: Product should equal exactly $77,000$
    \end{itemize}
\end{itemize}

\subsection*{D. Time Management}
\begin{itemize}[leftmargin=*]
    \item \textbf{Estimated Total Time}: 8-10 minutes
    \item \textbf{Time Checkpoints}:
    \begin{itemize}
        \item 2 min: Normalized problem and understood structure
        \item 5 min: Derived counting function $f(k)$
        \item 7 min: Found factorization $14 \times 50 \times 110$
        \item 8 min: Optimized and computed final answer
        \item 10 min: Verified answer
    \end{itemize}
    \item \textbf{Abandon Criterion}: If no progress on counting function after 5 minutes, flag and move on
\end{itemize}
\end{roadmapbox}

\begin{competitionbox}
\subsection*{A. Time Management Strategy}
\begin{itemize}[leftmargin=*]
    \item \textbf{Expected Time}: 8-10 minutes (this is P24, final problems deserve more time)
    \item \textbf{Decision Point}: After 5 minutes, assess progress:
    \begin{itemize}
        \item If counting function found: continue
        \item If stuck on counting: try answer verification approach
        \item If completely stuck: make educated guess and flag for review
    \end{itemize}
\end{itemize}

\subsection*{B. Problem Prioritization}
\begin{itemize}[leftmargin=*]
    \item \textbf{When to Attempt}: Only after completing P1-20 (easier problems)
    \item \textbf{Relative Difficulty}: This is harder than average P24 - requires insight
    \item \textbf{Domain Comfort}: High priority if strong in number theory and counting
\end{itemize}

\subsection*{C. Risk Management}
\begin{itemize}[leftmargin=*]
    \item \textbf{Confidence Threshold}:
    \begin{itemize}
        \item High (>85\%): Submit after standard verification
        \item Medium (70-85\%): Apply thorough verification (Level 5)
        \item Low (<70\%): Flag for review, return if time permits
    \end{itemize}
    \item \textbf{Error Recovery}: If arithmetic error found, restart calculation carefully
\end{itemize}

\subsection*{D. Abandonment Criteria}
\begin{itemize}[leftmargin=*]
    \item \textbf{Soft Abandon} (5 minutes): Move to P25, plan to return
    \item \textbf{Hard Abandon} (immediate):
    \begin{itemize}
        \item Can't derive counting function after multiple attempts
        \item Algebra becoming unmanageable
        \item Made same error multiple times
    \end{itemize}
    \item \textbf{Strategic Decision}: As P24, worth more time investment than earlier problems
\end{itemize}
\end{competitionbox}

\begin{keyinsightbox}
\textbf{Key Insight}: The critical breakthrough is recognizing that after normalizing by GCD, we can analyze each prime power independently. The counting function $f(k) = 12k^2 + 2$ follows from complementary counting: $(k+1)^4$ total choices minus cases missing either the maximum or minimum exponent. The factorization $77000 = 14 \times 50 \times 110 = f(1) \times f(2) \times f(3)$ immediately gives exponents, and minimizing requires assigning larger exponents to smaller primes.
\end{keyinsightbox}

\begin{warningbox}
\textbf{Warning}: The most common error is forgetting to multiply the final result by 77. After finding $N = 360$, students often submit this as the answer instead of $n = 77 \times 360 = 27,720$. Also, be careful with the exponent assignment—to minimize, put the largest exponent (3) on the smallest prime (2), not the other way around!
\end{warningbox}

\begin{tipbox}
\textbf{Pro Tip}: When deriving $f(k)$, you can either use the complementary counting formula or compute the first few values and recognize the pattern $f(k) = 12k^2 + 2$. For competition speed, computing $f(1) = 14, f(2) = 50, f(3) = 110$ and verifying $14 \times 50 \times 110 = 77000$ is faster than deriving the general formula. The pattern $12k^2 + 2$ emerges from plotting these points.
\end{tipbox}

\section*{Final Answer}
\[\boxed{\textbf{(D)}\ 27,720}\]

\section*{Detailed Solution}
\textbf{Step 1: Normalize by GCD}

Since $\gcd(a,b,c,d) = 77$, let $A = a/77, B = b/77, C = c/77, D = d/77$.

Then $\gcd(A,B,C,D) = 1$ and $N = n/77 = \operatorname{lcm}(A,B,C,D)$.

\textbf{Step 2: Analyze Prime Factorization}

Let $N = 2^{k_2} \times 3^{k_3} \times 5^{k_5} \times \cdots$

For each prime power $p^k$:
\begin{itemize}
    \item At least one of $(A,B,C,D)$ must have exponent $k$ (for LCM to be $N$)
    \item At least one must have exponent $0$ (for GCD to be $1$)
    \item All exponents must be in $[0,k]$
\end{itemize}

\textbf{Step 3: Counting Function}

Let $f(k)$ = number of ordered quadruplets $(m_1,m_2,m_3,m_4)$ with:
\begin{itemize}
    \item $0 \le m_i \le k$ for all $i$
    \item $\max(m_i) = k$ and $\min(m_i) = 0$
\end{itemize}

Using complementary counting:
\begin{align*}
f(k) &= \text{(all choices)} - \text{(missing max)} - \text{(missing min)} + \text{(missing both)} \\
&= (k+1)^4 - k^4 - k^4 + (k-1)^4 \\
&= 12k^2 + 2
\end{align*}

Verify: $f(1) = 14, f(2) = 50, f(3) = 110, f(4) = 194, \ldots$

\textbf{Step 4: Factor 77,000}

We need: $f(k_2) \times f(k_3) \times f(k_5) \times \cdots = 77,000$

Notice: $14 \times 50 \times 110 = 77,000$

Therefore: exponents are $(1, 2, 3)$

\textbf{Step 5: Minimize N}

To minimize $N = p_1^{e_1} \times p_2^{e_2} \times p_3^{e_3}$, assign larger exponents to smaller primes:

\begin{align*}
N &= 2^3 \times 3^2 \times 5^1 \\
&= 8 \times 9 \times 5 \\
&= 360
\end{align*}

Therefore: $n = 77 \times 360 = 27,720$

\section*{Confidence Assessment}
\begin{itemize}[leftmargin=*]
    \item \textbf{Confidence Level}: 90\% - Formula verified with multiple values, factorization confirmed, optimality checked
    \item \textbf{Verification Applied}: Level 5 (Thorough) - Forward check, formula verification, optimality analysis
    \item \textbf{Flag for Review}: No - high confidence after thorough verification
    \item \textbf{Time Spent}: ~10 minutes (appropriate for P24 complexity)
\end{itemize}

\end{document}